\documentclass[aspectratio=169]{beamer}
\usetheme{Bruno}
\usepackage{amsmath}
\usepackage{siunitx}
\usepackage[american,RPvoltages]{circuitikz}
\usepackage{tabularx}
\usepackage{tabu}
\usepackage[spanish]{babel}
\title{Electricidad I \\ \emph{Teoremas de circuitos}}

\usetikzlibrary{arrows, arrows.meta}
\author{
    Juan J. Rojas
}
\institute{Instituto Tecnológico de Costa Rica}
\date{\today}
\background{fig/background.jpg}
\renewcommand\tabularxcolumn[1]{m{#1}}% for vertical centering text in X column


\begin{document}

\maketitle
\sisetup{unit-math-rm=\mathrm,math-rm=\mathrm} % change sinitx font

\begin{frame}{Superposición}
    \begin{tabularx}{\linewidth}{X X}
        Pasos:
        \begin{enumerate}
            \item Se deja solamente una fuente independiente encendida. 
            \item Las otras fuentes independientes que no están siendo analizadas se apagan. Fuentes de voltaje se sustituyen por cortos; fuentes de corriente se sustituyen por abiertos.
        \end{enumerate}
        \vfill
        &
        \begin{enumerate}
        \setcounter{enumi}{2}
            \item Se calcula la contribución de la fuente analizada en el valor total de la corriente o el voltaje buscado.
            \item Se repiten los pasos 1, 2 y 3 para cada una de las fuentes independientes.
            \item Se suman las contribuciones de todas las fuentes independientes para determinar el valor total de la corriente o voltaje buscado. 
        \end{enumerate}
        \vfill
    \end{tabularx}
\end{frame}

\begin{frame}{Superposición}
        Solucionamos el ejemplo usado en supermallas. Buscamos $i_{R_6}$
        \vfill
        \centering
        \begin{circuitikz} [scale=0.65]\draw
            (0,0)
                to[V,l=$v_1$]
            (0,6)	
                to[R, l=$R_1$]
            (4,6)
                to[R,l=$R_3$]
            (8,6)
                to[R,l=$R_4$]
            (8,3)
                to[I,l=$i_1$,invert]
            (8,0)
                --
            (0,0)
            (4,0)
                to[R,l=$R_2$]
            (4,6)
            (8,6)
                to[R,l=$R_5$]
            (12,6)
                to[R,l=$R_6$, i=$i_{R_6}$]
            (12,0)
                --
            (8,0)
            ;
        \end{circuitikz}
\end{frame}

\begin{frame}{Superposición}
        1. Se deja solamente una fuente independiente encendida. 2. Las demás se apagan. 3. Se calcula la contribución de la fuente analizada en el valor total de la corriente o el voltaje buscado
        \vfill
        \centering
        \begin{circuitikz} [scale=0.65]\draw
            (0,0)
                to[V,l=$v_1$]
            (0,6)	
                to[R, l=$R_1$]
            (4,6)
                to[R,l=$R_3$]
            (8,6)
                to[open]
            (8,3)
                to[open]
            (8,0)
                --
            (0,0)
            (4,0)
                to[R,l=$R_2$]
            (4,6)
            (8,6)
                to[R,l=$R_5$]
            (12,6)
                to[R,l=$R_6$, i=$i_{R_{6v1}}$]
            (12,0)
                --
            (8,0)
            ;
        \end{circuitikz}
\end{frame}

\begin{frame}{Superposición}
        4. Se repiten los pasos 1, 2 y 3 para cada una de las fuentes independientes
        \vfill
        \centering
        \begin{circuitikz} [scale=0.65]\draw
            (0,0)
                to[short]
            (0,6)	
                to[R, l=$R_1$]
            (4,6)
                to[R,l=$R_3$]
            (8,6)
                to[R,l=$R_4$]
            (8,3)
                to[I,l=$i_1$,invert]
            (8,0)
                --
            (0,0)
            (4,0)
                to[R,l=$R_2$]
            (4,6)
            (8,6)
                to[R,l=$R_5$]
            (12,6)
                to[R,l=$R_6$, i=$i_{R_{6i1}}$]
            (12,0)
                --
            (8,0)
            ;
        \end{circuitikz}
\end{frame}

\begin{frame}{Superposición}
        5. Se suman las contribuciones de todas las fuentes independientes para determinar el valor total de la corriente o voltaje buscado.
        \vfill
        \centering
        \begin{circuitikz} [scale=0.65]\draw
            (0,0)
                to[V,l=$v_1$]
            (0,6)	
                to[R, l=$R_1$]
            (4,6)
                to[R,l=$R_3$]
            (8,6)
                to[R,l=$R_4$]
            (8,3)
                to[I,l=$i_1$,invert]
            (8,0)
                --
            (0,0)
            (4,0)
                to[R,l=$R_2$]
            (4,6)
            (8,6)
                to[R,l=$R_5$]
            (12,6)
                to[R,l=$R_6$, i=$i_{R_6}$]
            (12,0)
                --
            (8,0)
            ;
            \draw
            (-4,6) 
                node[]{$i_{R_6} = i_{R_{6v1}} + i_{R_{6i1}}$}
            ;
        \end{circuitikz}
\end{frame}

\begin{frame}{Transformación de fuentes}
        \vfill
        \centering
        \begin{circuitikz} [scale=0.8]\draw
            (4,0)
                to[short,o-]
            (0,0)
                to[V, v=$v_1$]
            (0,4)	
                to[R, l=$R_1$, -o]
            (4,4)
            (12,0)
                to[short,o-]
            (8,0)
                to[I, l=$i_1$]
            (8,4)
                to[short, -o]
            (12,4)
            (10,0)
                to[R, l_=$R_1$]
            (10,4)
            ;
            \draw[thick, >=triangle 45, <->] (4.5,2) -- (6.5,2);
        \end{circuitikz}
\end{frame}

\begin{frame}{Teorema de Thevenin}
        \vfill
        \centering
        \begin{circuitikz} [scale=0.8]\draw
            (1,-0.5)
                rectangle node[align=center] {\small Circuito\\\small lineal}
            (3,4.5)
            (3,4)	
                to[short,-o]
            (4,4)node[above]{$a$}
                to[short]
            (6,4)
                to[generic, l=\small{Carga}, i>^=$i$]
            (6,0)
                to[short, -o]
                node[below]{$b$}
            (4,0)
                to[short]
            (3,0)
            (4,4)
                to[open, v^=$v$]
            (4,0)
            (11,0)
                to[V,l=$v_{th}$]
            (11,4)
                to[R,l=$R_{th}$,-o]
                node[above]{$a$}
            (14,4)
                to[short]
            (16,4)
                to[generic, l=\small{Carga}, i>^=$i$]
            (16,0)
                to[short,-o]
                node[below]{$b$}
            (14,0)
                to[short]
            (11,0)
            (14,4)
                to[open, v^=$v$]
            (14,0)
            ;
            \draw[thick, >=triangle 45, ->] (8,2) -- (9.5,2);
        \end{circuitikz}
\end{frame}

\begin{frame}{Teorema de Thevenin}
        Se elimina la carga y podemos encontrar $v_{th}$ 
        \vfill
        \centering
        \begin{circuitikz} [scale=0.8]\draw
            (1,-0.5)
                rectangle node[align=center] {\small Circuito\\\small lineal}
            (3,4.5)
            (3,4)	
                to[short,-o]
            (4,4)node[above]{$a$}
            (4,0)node[below]{$b$}
                to[short,o-]
            (3,0)
            (4,4)
                to[open, v^=$v_{th}$]
            (4,0)
            (11,0)
                to[V,l=$v_{th}$]
            (11,4)
                to[R,l=$R_{th}$,-o]
            (14,4)node[above]{$a$}
                to[open, v^=$v_{th}$]
            (14,0)node[below]{$b$}
                to[short,o-]
            (11,0)
            ;
        \end{circuitikz}
\end{frame}

\begin{frame}{Teorema de Thevenin}
        Si no hay fuentes dependientes, se apagan las fuentes independientes y podemos encontrar $R_{th}$
        \vfill
        \centering
        \begin{circuitikz} [scale=0.8]\draw
            (1,-0.5)
                rectangle node[align=center] {\small Circuito\\\small lineal\\\small con\\\small fuentes\\\small indep.\\\small apagadas}
            (3,4.5)
            (3,4)	
                to[short,-o]
            (4,4)node[above]{$a$}
            (4,0)node[below]{$b$}
                to[short,o-]
            (3,0)
            (4,4)
                to[open, l=$R_{th}$]
            (4,0)
            (11,0)
                to[short]
            (11,4)
                to[R,l=$R_{th}$,-o]
            (14,4)node[above]{$a$}
                to[open, l=$R_{th}$]
            (14,0)node[below]{$b$}
                to[short,o-]
            (11,0)
            ;
            \draw[thick, >=triangle 45, <-] (4,1.6) -- (5,1.6);
            \draw[thick, >=triangle 45, <-] (14,1.6) -- (15,1.6);
        \end{circuitikz}
\end{frame}

\begin{frame}{Teorema de Thevenin}
        Si hay fuentes dependientes, se apagan las fuentes independientes y usamos una fuente de prueba para encontrar $R_{th}$
        \vfill
        \begin{tabularx}{\linewidth}{X X}
            \centering
            \begin{circuitikz} [scale=0.8]\draw
                (1,-0.5)
                    rectangle node[align=center] {\small Circuito\\\small lineal\\\small con\\\small fuentes\\\small indep.\\\small apagadas}
                (3,4.5)
                (3,4)	
                    to[short,-o]
                (4,4)node[above]{$a$}
                    to[short]
                (6,4)
                    to[V, l=$v_{o}$, i<^=$i_o$, invert]
                (6,0)
                    to[short, -o]
                (4,0)node[below]{$b$}
                    to[short,o-]
                (3,0)
                (6.5,4)node[right]{$R_{th} = \frac{v_o}{i_o}$}
                ;
            \end{circuitikz}
            &
            \centering
            \begin{circuitikz} [scale=0.8]\draw
                (1,-0.5)
                    rectangle node[align=center] {\small Circuito\\\small lineal\\\small con\\\small fuentes\\\small indep.\\\small apagadas}
                (3,4.5)
                (3,4)	
                    to[short,-o]
                (4,4)node[above]{$a$}
                    to[short]
                (6,4)
                    to[I, l=$i_{o}$, invert]
                (6,0)
                    to[short, -o]
                (4,0)node[below]{$b$}
                    to[short,o-]
                (3,0)
                (6.5,4)node[right]{$R_{th} = \frac{v_o}{i_o}$}
                (4,4)
                    to[open, v^=$v_o$]
                (4,0)
                ;
            \end{circuitikz}
        \end{tabularx}
\end{frame}

\begin{frame}{Teorema de Norton}
        \vfill
        \centering
        \begin{circuitikz} [scale=0.8]\draw
            (1,-0.5)
                rectangle node[align=center] {\small Circuito\\\small lineal}
            (3,4.5)
            (3,4)	
                to[short,-o]
            (4,4)node[above]{$a$}
                to[short]
            (6,4)
                to[generic, l=\small{Carga}, i>^=$i$]
            (6,0)
                to[short, -o]
                node[below]{$b$}
            (4,0)
                to[short]
            (3,0)
            (4,4)
                to[open, v^=$v$]
            (4,0)
            (11,0)
                to[I,l=$i_{N}$]
            (11,4)
                to[short,-o]
                node[above]{$a$}
            (14,4)
                to[short]
            (16,4)
                to[generic, l=\small{Carga}, i>^=$i$]
            (16,0)
                to[short,-o]
                node[below]{$b$}
            (14,0)
                to[short]
            (11,0)
            (13,4)
                to[R, l_=$R_N$]
            (13,0)
            (14,4)
                to[open, v^=$v$]
            (14,0)
            ;
            \draw[thick, >=triangle 45, ->] (8,2) -- (9.5,2);
        \end{circuitikz}
\end{frame}

\begin{frame}{Teorema de Norton}
        La carga se pone en cortocircuito y podemos encontrar $i_N$ 
        \vfill
        \centering
        \begin{circuitikz} [scale=0.8]\draw
            (1,-0.5)
                rectangle node[align=center] {\small Circuito\\\small lineal}
            (3,4.5)
            (3,4)	
                to[short,-o]
            (4,4)node[above]{$a$}
                to[short]
            (6,4)
                to[short, i>^=$i_N$]
            (6,0)
                to[short, -o]
                node[below]{$b$}
            (4,0)
                to[short,o-]
            (3,0)
            (11,0)
                to[I,l=$i_{N}$]
            (11,4)
                to[short,-o]
                node[above]{$a$}
            (14,4)
                to[short]
            (16,4)
                to[short, i>^=$i_N$]
            (16,0)
                to[short,-o]
                node[below]{$b$}
            (14,0)
                to[short]
            (11,0)
            (13,4)
                to[R, l_=$R_N$]
            (13,0)
            ;
        \end{circuitikz}
\end{frame}

\begin{frame}{Teorema de Norton}
    La resistencia equivalente de Norton es igual a la de Thevenin y se puede encontrar usando el mismo método.\\
    \begin{equation*}
        R_{th} = R_N
    \end{equation*}
    Luego se pueden relacionar por medio de la Ley de Ohm:
    \begin{equation*}
        i_N = \frac{v_{th}}{R_{th}}
    \end{equation*}
    \begin{equation*}
        v_{th} = i_N \cdot R_N
    \end{equation*}
\end{frame}

\begin{frame}{Teorema de máxima transferencia de potencia}
    Para transferir la máxima potencia a una carga, esta debe tener una resistencia igual a la resistencia de Thevenin.\\
    \vfill
    \centering
    \begin{circuitikz}[scale=0.8]\draw
        (0,0)
            to[V,l=$v_{th}$]
        (0,4)
            to[R,l=$R_{th}$,-o]
            node[above]{$a$}
        (3,4)
            to[short]
        (5,4)
            to[vR, mirror, l=$R_c$]
        (5,0)
            to[short,-o]
            node[below]{$b$}
        (3,0)
            to[short]
        (0,0)
        (3,4)
            to[open, v^=$v$]
        (3,0)
        (8,4)node[right]{$p=\left(\frac{v_{th}}{R_{th} + R_c}\right)^2 \cdot R_c$}
        (8,2)node[right]{$p_{max} \rightarrow R_c = R_{th}$}
        (8,0)node[right]{$p_{max}=\frac{v_{th}^2}{4\cdot R_{th}}$}
        ;
    \end{circuitikz}
\end{frame}

\begin{frame}{Puente de Wheatstone}
        \begin{tabularx}{\linewidth}{X X}
        \begin{equation*}
            v_a = v_1
        \end{equation*}
        \begin{multline*}
             -\left(\frac{1}{R_1}+\frac{1}{R_2}\right) \cdot v_a + \left(\frac{1}{R_2} + \frac{1}{R_3}\right) \cdot v_b \\+ \left(\frac{1}{R_1} + \frac{1}{R_4}\right) \cdot v_c= 0
        \end{multline*}
        \vspace{0.2cm}
        \begin{equation*}
            v_b - v_c = v_2
        \end{equation*}
        &
        \centering
        \begin{circuitikz} [scale=0.8]\draw
            (0,0)
                to[V,l=$v_1$]
                node[left]{$v_a$}
            (0,4)
                to[short]
            (0,6)
                to[R,l=$R_1$]
                %node[ground]{}
            (6,6)
                to[short]
                node[right]{$v_c$}
            (6,4)
                to[R,l=$R_4$]
            (6,0)
                to[short]
                node[ground]{}
            (3,0)
                to[short]
            (0,0)
            (0,4)
                 to[R,l=$R_2$]
                 node[above]{$v_b$}
            (3,4)
                to[R,l=$R_3$]
            (3,0)
            (6,4)
                to[V,l_=$v_2$]
            (3,4)
            ;
        \end{circuitikz}
        \vfill
        \end{tabularx}
\end{frame}

\end{document}